\section{Height/regulator column}
\label{sec:height-regulator}

The height/regulator column keeps the Mordell--Weil lattice and its
N\'eron--Tate pairing~\citep{Neron1965,Silverman2009} before any comparison
with the global determinant line.  Its readout is the regulator, viewed as
the coefficient of a determinant.

\begin{definition}[Real height-regulator map]\label{def:apparatus:height-regulator-map}
Let
\[
  \Lambda_E=E(\Q)/E(\Q)_{\tors},
  \qquad
  V_E=\Lambda_E\otimes_{\Z}\R,
\]
and let $\langle\cdot,\cdot\rangle_{\mathrm{NT}}$ be the N\'eron--Tate
pairing on $V_E$, with associated map
\[
  h_{\mathrm{NT}}:V_E\longrightarrow V_E^*,
  \qquad
  P\longmapsto \langle P,\cdot\rangle_{\mathrm{NT}}.
\]
The real height-regulator map sends the triple
$(V_E,\Lambda_E,\langle\cdot,\cdot\rangle_{\mathrm{NT}})$ to
\[
  \rho_{\mathrm{ht/reg}}^{\R}
  =\det(h_{\mathrm{NT}})\in\det(V_E^*)\otimes\det(V_E)^*.
\]
If $(P_1,\ldots,P_r)$ is a $\Z$-basis of $\Lambda_E$ and
$H_E=\bigl(\langle P_i,P_j\rangle_{\mathrm{NT}}\bigr)_{1\le i,j\le r}$ is
the Gram matrix, the coefficient of this element in the induced basis is
$\det H_E=\Reg(E/\Q)$, with the empty determinant equal to $1$ when $r=0$.
The coefficient does not depend on the basis: a change of basis by
$A\in\mathrm{GL}_r(\Z)$ replaces $H_E$ by $A^{\mathsf T}H_EA$, and
$\det(A^{\mathsf T}H_EA)=\det(A)^2\det H_E=\det H_E$ because
$\det A=\pm1$.  (A change of basis of finite index $m>1$ multiplies the
coefficient by $m^2$, so integrality of the basis matters.)
\end{definition}

The adequacy test for this column asks whether the full height matrix is
enough to recover the first-order variation of $\log\Reg$.  It is, for a
simple reason: the readout is a linear function of the variation, and the
full source sees every coordinate of the variation.

\begin{theorem}[Height tangent Schur collapse]\label{thm:apparatus:height-schur-collapse}
Let $r\ge1$ and let $H$ be a positive definite real symmetric $r\times r$
matrix.  Identify a symmetric variation $\delta H$ with its coordinate
vector $x=\operatorname{vech}(\delta H)\in\R^n$, $n=r(r+1)/2$, listing the
entries $\delta h_{ij}$ with $i\le j$.  Then
\[
  D_H(\delta H)
  =
  d(\log\det)_H(\delta H)
  =
  \tr(H^{-1}\delta H)
  =
  d_Hx,
\]
where the row $d_H$ has entry $(H^{-1})_{ii}$ at each diagonal
coordinate and $2(H^{-1})_{ij}$ at each off-diagonal coordinate $i<j$.
Take the audit energy $C=I_n$ on these coordinates and the full source
$L_{\mathrm{full}}(x)=x$.  Then
\[
  \KDD=\|d_H\|^2,
  \qquad
  \XiC=d_Hd_H^{\mathsf T}-d_H\,I_n\,d_H^{\mathsf T}=0 .
\]
For $r=0$ the tangent space is zero, $n=0$, and $\KDD=0$: the rank-zero
row is vacuous.%
\footnote{Lean proves the algebraic core: for every row $d$ of every
length over any commutative ring, the identity source gives Schur
residual zero, and the identity matrix is the unique matrix satisfying
the four Penrose equations at the identity.  For $r=2$ it also proves
that the displayed row computes $\tr(H^{-1}\delta H)$ when
$\det H\neq0$.  The calculus step identifying the row with the derivative
of $\log\det$, and the general-rank row, are proved here on paper.}
\end{theorem}

\begin{proof}
Since $H$ is invertible, $\det(H+t\,\delta H)=\det H\cdot\det(I+tA)$ with
$A=H^{-1}\delta H$.  In the permutation expansion of $\det(I+tA)$, the
identity permutation contributes $1+t\tr A+O(t^2)$, and every other
permutation moves at least two indices and contributes $O(t^2)$.  Since
$\det H>0$, the chain rule gives $d(\log\det)_H(\delta H)=\tr(H^{-1}\delta
H)$.  Writing $G=H^{-1}$, which is symmetric,
\[
  \tr(G\,\delta H)=\sum_{i,j}G_{ij}\,\delta h_{ji}
  =\sum_i G_{ii}\,\delta h_{ii}+\sum_{i<j}2G_{ij}\,\delta h_{ij},
\]
because the off-diagonal coordinate $\delta h_{ij}$ occurs twice in the
symmetric matrix $\delta H$.  This is $d_Hx$.

For the Schur residual, the source block is $\KLL=I_n$, whose
Moore--Penrose inverse is $I_n$, and the cross blocks are $\KDL=d_H$ and
$K_{LD}=d_H^{\mathsf T}$.  Hence
$\KDL\Kdagger K_{LD}=d_H\,I_n\,d_H^{\mathsf T}=d_Hd_H^{\mathsf T}=\KDD$,
and $\XiC=0$.  If $r=0$, the row $d_H$ is empty and the empty sum
$d_Hd_H^{\mathsf T}$ is $0$.
\end{proof}

For example, when $r=2$ and $H=\bigl(\begin{smallmatrix}a&b\\
b&c\end{smallmatrix}\bigr)$ with $\Delta=ac-b^2>0$, the row in the
coordinates $(\delta h_{11},\delta h_{12},\delta h_{22})$ is
\[
  d_H=\Bigl(\frac{c}{\Delta},\,-\frac{2b}{\Delta},\,\frac{a}{\Delta}\Bigr),
\]
as one also sees from $\det(H+t\,\delta H)=\Delta+t\,(c\,\delta h_{11}-2b\,
\delta h_{12}+a\,\delta h_{22})+O(t^2)$.  When $r=1$, $d_H=1/\Reg(E/\Q)$
and $\KDD=\Reg(E/\Q)^{-2}$; for the curve $37a1$, the regulator
$0.0511114\ldots$ recorded in the LMFDB~\citep{LMFDB} gives
$\KDD=382.79\ldots$.  The choice $C=I_n$ refers to the chosen
$\operatorname{vech}$ coordinates; it is not invariant under changes of
lattice basis, and only the vanishing of $\XiC$, not the value of $\KDD$,
is used later.

\begin{obligation}[Height/regulator transport]\label{obl:apparatus:rho-ht-reg-transport}
The Bloch--Kato height/regulator component requires a comparison
\[
  \beta_{\mathrm{BK,ht}}:
  \det(V_E^*)\otimes\det(V_E)^*
  \longrightarrow
  \Delta_f(M_E)_{\mathrm{ht/reg}}
\]
identifying the N\'eron--Tate determinant with the Beilinson--Deligne
regulator normalization~\citep{Beilinson1985}.  It must fix
the coefficient and lattice conventions, the archimedean
Deligne-cohomology normalization, and the orientation and
trivialization, and it must be compatible with the finite and
determinant-assembly columns without absorbing them.  This comparison is
external arithmetic input.  It is not derived from the finite calculation
above, which only shows that the full height matrix is adequate for the
first-order regulator readout.
\end{obligation}

\begin{warning}[Rank-zero carriers]\label{warn:apparatus:571a1-carrier}
For a curve of rank $0$ the height/regulator column is empty: there is no
nonzero height variation, and the regulator is the empty determinant $1$.
The curve $571a1$, which reappears in \Cref{sec:finite-source} because
its Tate--Shafarevich group has nontrivial $2$-primary part, is such a
curve: a certified $2$-descent bounds its rank above and below by $0$
(details in \Cref{sec:finite-source}).  It should not be confused with
the rank-two curve $571b1$ of the same conductor.  Rank-two
height/regulator examples in this paper use $389a1$.
\end{warning}
